\chapter*{Tóm tắt}

Các bệnh lý và chấn thương xương ngày càng phổ biến, làm gia tăng nhu cầu ứng dụng trí tuệ nhân tạo trong phân tích ảnh X-quang. Tuy nhiên, dữ liệu y khoa thường hạn chế về số lượng mẫu gán nhãn, mất cân bằng lớp và tồn tại bệnh ngoài phân phối, giảm khả năng tổng quát của mô hình khi triển khai. Đề tài nhằm xây dựng hệ thống phân loại đa lớp ảnh X-quang xương theo hướng đa phương thức, khai thác đồng thời hình ảnh và văn bản lâm sàng để nâng cao hiệu quả dự đoán. Hệ thống được phát triển trên nền tảng BioMedCLIP với hai giai đoạn huấn luyện: thích nghi miền bằng LoRA kết hợp hàm mất mát khớp ngữ nghĩa để cải thiện biểu diễn đặc trưng, sau đó hợp nhất đặc trưng ảnh và văn bản bằng cơ chế chú ý hai chiều. Mạng nguyên mẫu kết hợp Cross-Entropy có trọng số lớp được sử dụng nhằm giảm ảnh hưởng mất cân bằng lớp, hỗ trợ các lớp ít mẫu và phát hiện dữ liệu ngoài phân phối dựa trên khoảng cách mahalanobis. Phương pháp được đánh giá thông qua thí nghiệm thành phần, so sánh với mô hình nền và kiểm thử trên bộ dữ liệu tại Bệnh viện Chấn thương Chỉnh hình. Kết quả cho thấy phương pháp đề xuất cải thiện hiệu quả phân loại đa lớp và khả năng tổng quát, đồng thời hỗ trợ nhận diện mẫu ngoài phân phối.